% !TeX root = ../../../main.tex
\subsection{Data Records}
\label{ch05:sec:data-records}

\ChFiveRebuttal{
The PMOA-TTS dataset is deposited on Hugging Face \cite{shahriar_noroozizadeh_2026} and is publicly available at \url{https://doi.org/10.57967/hf/10546}.
}
\ChFiveRebuttal{
PMOA-TTS distributes derived structured annotations and source identifiers rather than copies of source-article text, figures, images, or supplementary files. The released annotations and metadata are licensed under CC BY-NC-SA 4.0 (\url{https://creativecommons.org/licenses/by-nc-sa/4.0/}); reuse of the underlying source articles remains subject to the article-level terms displayed by PMC (\url{https://pmc.ncbi.nlm.nih.gov/tools/openftlist/}).
}
The dataset card lists task tags (\texttt{Text Classification}), modality (\texttt{Text}), format support (\texttt{Parquet via auto-conversion}), language (\texttt{English}), and size (``\texttt{100K--1M}'').
Data record schema with an example JSON entry is shown in Figure~\ref{ch05:fig:record_schema}.
% \ChFiveRebuttal{
% The demographic fields are LLM-extracted or LLM-mapped descriptors from the case-report narrative. They should not be interpreted as independently verified biological sex at birth, self-identified race or ethnicity, or population-representative demographic measurements. The model returns Not Specified when no value is identified; ethnicity is Not Specified in 89\% of reports.
% }

\noindent PMOA-TTS has the following dataset splits:

% !TeX root = ../../../main.tex
\begin{figure}[htbp]
\begin{ChFiveFloatBody}

\centering

\begin{minipage}[t]{1\textwidth}
\centering
\footnotesize
\begin{adjustbox}{max width=\linewidth}\begin{tabular}{@{} l l >{\raggedright\arraybackslash}p{4cm} @{}}
\toprule
\textbf{Field} & \textbf{Type} & \textbf{Description} \\
\midrule
\texttt{pmc\_id} & str &
Folder-level ID (PMC000xxxxxx--PMC011xxxxxx). \\

\texttt{case\_report\_id} & str &
Filename of the case (e.g., PMC6417290). \\

\texttt{textual\_timeseries} & List[Dict] &
Sequence of \{"event": str, "time": int\}. \\

\texttt{demographics} & Dict &
\{"age": str/"Not Specified", "sex": str, "ethnicity": str\}. \\

\texttt{diagnoses} & List[str] &
List of diagnosis terms extracted per case. \\

\bottomrule
\end{tabular}\end{adjustbox}
\end{minipage}
\hfill

\begin{minipage}[t]{\linewidth}
\centering
\begin{tcolorbox}[
  colback=gray!5,
  colframe=gray!70,
  title=Example JSON entry,
  boxrule=1pt,
  enhanced,
  %breakable,
  fontupper=\footnotesize
]
\begin{verbatim}
"pmc_id": "PMC006xxxxxx",
"case_report_id": "PMC6417290",

"textual_timeseries": [
  {"event": "56 years old", "time": 0},
  {"event": "male", "time": 0},
  {"event": "HIV-positive", "time": 0},
  {"event": "admitted to the hospital", "time": 0},
  {"event": "knee arthralgia", "time": 0},
  ...
  {"event": "postural headache", "time": 48},
  {"event": "neurological examination", "time": 120},
  {"event": "persistence of headache", "time": 144},
  {"event": "brain computed tomography scan",
   "time": 144},
  ...
  {"event": "symptom free", "time": 4320},
  {"event": "CT scan", "time": 4320},
  {"event": "no brain shift", "time": 4320}
],
"diagnoses": [
  "HIV",
  "Knee arthralgia",
  "Subdural hematoma",
  "Postdural puncture headache (PDPH)"
],
"demographics": {
  "age": 56,
  "sex": "Male",
  "ethnicity": "Not Specified"
}
\end{verbatim}
\end{tcolorbox}
\end{minipage}





\end{ChFiveFloatBody}
\caption{\textbf{top:} PMOA-TTS record schema. \textbf{bottom:} compact example matching the schema above (drawn from our HuggingFace repository (formatted for readability).}\ChFiveFloatLabel{ch05:fig:record_schema}
\end{figure}




\begin{itemize}
    \item \texttt{train\_DSR1/L33}: The full TTS release contains \ChFiveRebuttal{124{,}349 records} with valid event-time rows.
The \texttt{train\_DSR1} and \texttt{train\_L33} splits contain annotations generated using \texttt{DeepSeek-R1} and \texttt{Llama 3.3 70B}, respectively.
Here, a split denotes a version of the dataset defined by the annotation model used.
These splits serve as the default releases and contain the primary annotations used in downstream analyses.
    
    \item {\raggedright \ChFiveRebuttal{\texttt{case\_study\_100} contains the TTS-bearing portion of the 100-document case-identification evaluation. Of the 100 candidate documents reviewed, 88 were confirmed to contain a single-patient case report and therefore have TTS annotations; the remaining 12 did not contain a case report and consequently have no patient timeline in this split.}\par}
    
    \item \texttt{case\_study\_n200}: A benchmark of 200 manually reviewed case reports used for qualitative and quantitative evaluation of LLM-generated annotations.
Each model is evaluated across five independent runs on this subset to assess run-to-run variability in annotation consistency.

\end{itemize}

The data repository also has an additional data split \texttt{TTS\_n200\_additional\_models}, which bundles the single-expert manual timelines (the reference set) together with parallel LLM-generated textual time series for additional models, including open-weight models such as \texttt{DeepSeek-V3}, \texttt{DeepSeek-R1 quantized}, and \texttt{DeepSeek-V3-0324}, and proprietary OpenAI models such as \texttt{GPT5}, \texttt{O1}, \texttt{O3}, and \texttt{O4-mini}.
This split is designed for lightweight benchmarking and cross-model comparison on a fixed 200-case cohort without re-annotating the full corpus.
For the full \ChFiveRebuttal{124{,}349} release, we provide textual time series produced by \texttt{DeepSeek-R1} and \texttt{Llama 3.3 70B}; the additional models are released only on the 200-case subset due to resource constraints.
For completeness, both \texttt{Llama 3.3 70B} and \texttt{DeepSeek-R1} were also executed in five independent runs on the \texttt{case\_study\_n200} subset, and these annotated outputs are included in the dataset.
\ChFiveRebuttal{Additional details regarding exact model versions, access dates, decoding settings, context limits, seeds, and the inference framework are available in Supplementary Material Section  \ref{ch05:sec:supp-model-configuration}}.

